\documentclass[11pt]{article}
% ---------- Pakete ----------
\usepackage[utf8]{inputenc}
\usepackage[T1]{fontenc}
\usepackage[ngerman,provide=*]{babel}
\usepackage[a4paper,top=0.0cm,bottom=0.55cm,left=1.3cm,right=1.3cm]{geometry}
\usepackage{helvet}
\renewcommand{\familydefault}{\sfdefault}
\usepackage{xcolor}
\usepackage{amssymb}
\usepackage{array}
\usepackage{tikz}
\usetikzlibrary{positioning}
\usepackage[most]{tcolorbox}
\usepackage{enumitem}
\usepackage{pifont}
\usepackage{setspace}
\setstretch{1.03}
\setlist{nosep,leftmargin=1.1em}
\pagestyle{empty}
\setlength{\parindent}{0pt}
% ---------- Farben ----------
\definecolor{egreen}{RGB}{27,128,77}
\definecolor{egreendark}{RGB}{18,92,55}
\definecolor{egreenlight}{RGB}{233,245,238}
\definecolor{amber}{RGB}{214,158,46}
\definecolor{boxgray}{RGB}{110,110,110}
\definecolor{greenverlauf}{RGB}{25 195 125}
\definecolor{turkisverlauf}{RGB}{14 140 158}
\definecolor{cardgray}{RGB}{244,246,245}
\definecolor{turkislight}{RGB}{230,242,245}
% ---------- Hilfsbefehle ----------
\newsavebox\imgbox
% Grüne Abschnitts-Überschrift (wie im Mockup)
\newcommand{\secheading}[1]{%
  \par{\large\bfseries\color{egreendark}#1}\par\vspace{4pt}}
% Icon-Chip: kleines gerundetes Quadrat mit Symbol
\newcommand{\chip}[1]{%
  \tikz[baseline=-4pt]\node[rounded corners=2.5pt, fill=white,
        inner sep=3.5pt]{\color{turkisverlauf}#1};}
% Vorteils-Box (in tcbraster verwenden)
\newcommand{\vorteil}[3]{%
  \begin{tcolorbox}[colback=turkislight, colframe=turkislight, boxrule=0pt, arc=6pt,
                    left=5pt, right=5pt, top=4pt, bottom=4pt, valign=top]
    \raggedright
    \chip{#1}~{\bfseries #2}\par\vspace{2.5pt}
    {\small #3}
  \end{tcolorbox}}
% Nummern-Kreis für Schritte
\newcommand{\stepnum}[1]{%
  \tikz[baseline=-4pt]\node[circle, fill=egreendark, text=white,
        inner sep=1.5pt, minimum size=15pt, font=\bfseries\small]{#1};}
\begin{document}
% =====================================================
% KOPF (volle Breite, Farbverlauf)
% =====================================================
\begin{tcolorbox}[enhanced, interior style={left color=greenverlauf, right color=turkisverlauf},
                  colframe=egreendark, sharp corners, boxrule=0pt,
                  left=10pt, right=10pt, top=8pt, bottom=8pt,
                  grow to left by=1.3cm, grow to right by=1.3cm]
\begin{minipage}[c]{2.6cm}
  \includegraphics[width=2.2cm]{icon-white.png}
\end{minipage}\hfill
\begin{minipage}[c]{\dimexpr\linewidth-3.0cm\relax}
  \color{white}
  \vspace{5pt}
  {\fontsize{21}{24}\selectfont\bfseries Sonnenstrom aus deiner Nachbarschaft}\\[6pt]
  {\normalsize Der \textbf{Strompool Feldkirchen Süd-West} teilt Sonnenstrom in der
  Region \\ fair, günstig und ohne Risiko. Mach mit, mit oder ohne eigene PV-Anlage.}
\end{minipage}
\end{tcolorbox}
\vspace{9pt}
% =====================================================
% WAS IST EINE EEG? (Text links, Diagramm rechts)
% =====================================================
\begin{minipage}[t]{0.55\linewidth}
\secheading{Was ist eine Erneuerbare Energiegemeinschaft?}
Eine Erneuerbare Energiegemeinschaft (\textbf{EEG}) ist ein Verein von Nachbarn.
Haushalte und Betriebe aus der Region teilen selbst erzeugten Sonnenstrom.
Der Strom fließt über das \textbf{bestehende Stromnetz} -- ganz ohne neue
Leitungen.\par\vspace{6pt}
\begin{tcolorbox}[enhanced, colback=egreenlight, colframe=egreenlight, boxrule=0pt, arc=3pt,
                  borderline west={2.5pt}{0pt}{egreen},
                  left=8pt, right=6pt, top=5pt, bottom=5pt]
\small \textbf{Wie ein Gemeinschaftsgarten -- nur für Strom:} \\Wer zu viel
erntet, gibt günstig an die Nachbarn ab.
\end{tcolorbox}
\vspace{6pt}
Möglich macht das ein Gesetz: das Erneuerbaren-Ausbau-Gesetz (seit 2021).
Mitmachen kann jeder -- \textbf{mit PV-Anlage} (verkauft den Überschuss)
oder \textbf{ohne} (bezieht günstigen Regionalstrom).
\end{minipage}\hfill
\begin{minipage}[t]{0.41\linewidth}
\vspace{2pt}

\centering

\sbox\imgbox{\includegraphics[width=\linewidth]{EEG Grafik.png}}%
\begin{tikzpicture}
  \clip[rounded corners=6pt] (0,0) rectangle (\wd\imgbox,\ht\imgbox);
  \node[anchor=south west, inner sep=0pt] at (0,0) {\usebox\imgbox};
\end{tikzpicture}

{\footnotesize\color{boxgray} Überschuss-Strom wandert \\ übers normale Netz zu den Nachbarn.}

\end{minipage}
\vspace{9pt}
% =====================================================
% TARIF-KACHELN
% =====================================================
\begin{tcbraster}[raster columns=3, raster equal height, raster column skip=9pt,
                  raster before skip=0pt, raster after skip=0pt]
\begin{tcolorbox}[colback=greenverlauf, colframe=greenverlauf, boxrule=0pt, arc=6pt,
                  left=4pt, right=4pt, top=6pt, bottom=6pt, valign=top]
\centering\color{white}{\huge\bfseries 12 ct}{\small\bfseries /kWh}\\[2pt]
{\small Strombezug aus der \\Gemeinschaft}
\end{tcolorbox}
\begin{tcolorbox}[colback=turkisverlauf, colframe=turkisverlauf, boxrule=0pt, arc=6pt,
                  left=4pt, right=4pt, top=6pt, bottom=6pt, valign=top]
\centering\color{white}{\huge\bfseries 8 ct}{\small\bfseries /kWh}\\[2pt]
{\small Vergütung für deinen PV-Überschuss}
\end{tcolorbox}
\begin{tcolorbox}[colback=white, colframe=egreen, boxrule=1.2pt, arc=6pt,
                  left=4pt, right=4pt, top=6pt, bottom=6pt, valign=top]
\centering\color{egreendark}{\huge\bfseries 2 \texteuro}{\small\bfseries /Monat}\\[2pt]
{\small Mitgliedsbeitrag \\ kein Eintrittsgeld}
\end{tcolorbox}
\end{tcbraster}
\vspace{9pt}
% =====================================================
% VORTEILE
% =====================================================
\secheading{Deine Vorteile}
\vspace{4pt}

\begin{tcbraster}[raster columns=3, raster equal height, raster column skip=7pt,
                  raster row skip=7pt, raster before skip=0pt, raster after skip=0pt]
\vorteil{\bfseries\texteuro}{Günstiger Strom}{12 ct/kWh -- meist deutlich unter dem
normalen Liefertarif.}
\vorteil{\ding{72}}{Faire Vergütung}{8 ct/kWh für deinen \\ PV-Überschuss... 
\\oft mehr als üblich.}
\vorteil{\ding{34}}{Weniger Abgaben}{Rund 28\,\% weniger Netzgebühr
und keine Stromsteuer auf EEG-Strom.}
\vorteil{\ding{170}}{Regional \& transparent}{Verein ohne Gewinnabsicht. 
\\Die Wertschöpfung bleibt in Feldkirchen.}
\vorteil{\ding{96}}{Klimaschutz}{100\,\% Sonnenstrom \\aus der Region.}
\vorteil{\ding{52}}{Kein Risiko}{Dein Stromvertrag bleibt aufrecht.
Ausstieg jederzeit möglich.}
\end{tcbraster}
\vspace{9pt}
% =====================================================
% SO MACHST DU MIT (4 Karten)
% =====================================================
\secheading{So machst du mit}
\vspace{4pt}

\begin{tcbraster}[raster columns=4, raster equal height, raster column skip=7pt,
                  raster before skip=0pt, raster after skip=0pt]
\begin{tcolorbox}[colback=cardgray, colframe=cardgray, boxrule=0pt, arc=6pt,
                  left=6pt, right=6pt, top=4pt, bottom=4pt, valign=top]
\raggedright\stepnum{1}\par\vspace{4pt}
{\small \textbf{Online beitreten} \\auf \texttt{stromfueralle.at} unter „Beitreten“.}
\end{tcolorbox}
\begin{tcolorbox}[colback=cardgray, colframe=cardgray, boxrule=0pt, arc=6pt,
                  left=6pt, right=6pt, top=4pt, bottom=4pt, valign=top]
\raggedright\stepnum{2}\par\vspace{4pt}
{\small \textbf{Zustimmung beim Netzbetreiber} \\(Kärnten Netz)}
\end{tcolorbox}
\begin{tcolorbox}[colback=cardgray, colframe=cardgray, boxrule=0pt, arc=6pt,
                  left=6pt, right=6pt, top=4pt, bottom=4pt, valign=top]
\raggedright\stepnum{3}\par\vspace{4pt}
{\small \textbf{Smart Meter} \\(digitaler Stromzähler) teilt den Strom automatisch zu.}
\end{tcolorbox}
\begin{tcolorbox}[colback=cardgray, colframe=cardgray, boxrule=0pt, arc=6pt,
                  left=6pt, right=6pt, top=4pt, bottom=4pt, valign=top]
\raggedright\stepnum{4}\par\vspace{4pt}
{\small \textbf{Abrechnung} \\kommt bequem von \\der EEG pro Quartal oder pro Jahr.}
\end{tcolorbox}
\end{tcbraster}
\vspace{5pt}
{\small \textbf{Wichtig:} An Leitung, Zähler und Versorgungssicherheit ändert sich nichts.}
\vspace{9pt}
% =====================================================
% FAQ (3 Spalten mit Akzentlinie)
% =====================================================
\secheading{Häufige Fragen}
\vspace{4pt}

\begin{tcbraster}[raster columns=3, raster equal height, raster column skip=9pt,
                  raster before skip=0pt, raster after skip=0pt]
\begin{tcolorbox}[enhanced, colback=white, colframe=white, boxrule=0pt, arc=0pt,
                  borderline west={2.5pt}{0pt}{turkisverlauf},
                  left=8pt, right=2pt, top=1pt, bottom=1pt, valign=top]
\raggedright\small
\textbf{„Brauche ich eine eigene PV-Anlage?“}\par\vspace{2pt}
Nein. Auch reine Strom-Verbraucher sind herzlich willkommen.
\end{tcolorbox}
\begin{tcolorbox}[enhanced, colback=white, colframe=white, boxrule=0pt, arc=0pt,
                  borderline west={2.5pt}{0pt}{turkisverlauf},
                  left=8pt, right=2pt, top=1pt, bottom=1pt, valign=top]
\raggedright\small
\textbf{„Was, wenn die Sonne nicht scheint?“}\par\vspace{2pt}
Dein normaler Stromvertrag bleibt bestehen -- es fehlt nie Strom.
\end{tcolorbox}
\begin{tcolorbox}[enhanced, colback=white, colframe=white, boxrule=0pt, arc=0pt,
                  borderline west={2.5pt}{0pt}{turkisverlauf},
                  left=8pt, right=2pt, top=1pt, bottom=1pt, valign=top]
\raggedright\small
\textbf{„Was kostet der Beitritt?“}\par\vspace{2pt}
2 \texteuro/Monat. Kein Eintrittsgeld, keine versteckten
Kosten, jederzeit kündbar.
\end{tcolorbox}
\end{tcbraster}
\vfill
% =====================================================
% CTA + FUSS
% =====================================================
\begin{tcolorbox}[enhanced, interior style={left color=greenverlauf, right color=turkisverlauf},
                  colframe=egreendark, sharp corners, boxrule=0pt,
                  left=8pt, right=10pt, top=3pt, bottom=3pt]
\color{white}
{\large\bfseries Jetzt Mitglied werden:} \quad
{\large\texttt{stromfueralle.at/rc108175/beitreten}}\\[4pt]
\footnotesize EEG Strompool Feldkirchen Süd-West \;$\cdot$\; ZVR 1778816746 \;$\cdot$\;
RC 108175 \;$\cdot$\; Sitz: Feldkirchen in Kärnten\\
Kontakt: Patrick Ropper, \mbox{office@stromfueralle.at}, Tel.\ +43\,650\,6044812
\;$\cdot$\; Angaben ohne Gewähr, Stand Juli 2026
\end{tcolorbox}
\end{document}